\subsection{Typy výstupů a iterativní ladění}
\label{subsec:prompts-output}

Výstup jazykového modelu při analýze studentského kódu může mít několik podob v závislosti na tom, jaký typ zpětné vazby je pro učitele a studenty nejužitečnější.

\subsubsection*{Typy výstupů}
\label{subsubsec:prompts-output-types}

\begin{itemize}
    \item \textbf{Komentáře k řádkům} představují nejpodrobnější formu zpětné vazby.
    Každý nalezený problém je vázán na konkrétní řádek zdrojového kódu a obsahuje popis problému spolu s návrhem opravy.
    Tento formát je pro studenty nejpřínosnější, protože přesně ukazuje, kde v kódu se problém nachází, a nabízí konkrétní vysvětlení.

    \item \textbf{Souhrnné hodnocení} poskytuje celkový pohled na kvalitu odevzdaného řešení.
    Může zahrnovat informaci o tom, zda řešení splňuje požadavky zadání, jaká je jeho celková úroveň a jaké hlavní oblasti vyžadují zlepšení.
    Tento typ výstupu je užitečný především pro učitele jako rychlý přehled před detailním hodnocením.
\end{itemize}

Pro integraci do systému Kelvin byl zvolen formát, který kombinuje oba výše zmíněné typy výstupy v podobně JSON souboru.
Tento formát umožňuje programové zpracování na straně backendu a zároveň tak jednoduše předává tyto informace na frontend.

\subsubsection*{Struktura výstupního formátu}
\label{subsubsec:prompts-output-structure}

Výstup modelu je definován JSON schématem, které je součástí systémového promptu a zajišťuje, že model vrací data v přesně specifikované struktuře.
Toto schéma je předáno modelu prostřednictvím parametru \texttt{response\_format}, čímž je vynucen strukturovaný výstup bez nutnosti dodatečného zpracování volného textu.

Výsledek analýzy je reprezentován objektem \texttt{ReviewResult}, který obsahuje dva prvky: textové shrnutí celkové kvality kódu a seznam konkrétních nalezených problémů.
Shrnutí (\texttt{summary}) poskytuje učiteli rychlý přehled o architektuře, čitelnosti a udržovatelnosti analyzovaného kódu, aniž by opakovalo jednotlivé nálezy.

Každý problém v seznamu (\texttt{issues}) je reprezentován objektem \texttt{ReviewIssue} se čtyřmi povinnými atributy:

\begin{itemize}
    \item \textbf{file} určuje název souboru, ve kterém byl problém identifikován.

    \item \textbf{severity} vyjadřuje závažnost nálezu pomocí čtyřstupňové škály: \texttt{critical} pro chyby, které zásadně narušují funkčnost programu, \texttt{high} pro závažné problémy ovlivňující správnost, \texttt{medium} pro nedostatky v kvalitě kódu a \texttt{low} pro drobné stylistické připomínky.

    \item \textbf{line} obsahuje číslo řádku, na kterém se problém nachází.
    Schéma explicitně instruuje model, aby uváděl řádek prvního tokenu přímo zodpovědného za daný problém, nikoli uzavírací závorku, prázdný řádek nebo komentář.
    Pokud se problém týká více řádků, uvede se nejranější příčinný řádek a ostatní jsou zmíněny v textovém vysvětlení.

    \item \textbf{explanation} obsahuje stručný popis problému: co je chybou, proč je to problém a jak jej opravit.
    Schéma vyžaduje, aby model odkazoval na konkrétní názvy proměnných a funkcí a vyhýbal se obecným frázím bez informační hodnoty.
\end{itemize}

\endinput